\chapter{Data loading}
\label{ch:data-loading}

\begin{quote}
\emph{``You can't analyze what you can't load.''}
\end{quote}

% ============================================================
\section{Loading CSV files}

CSV (comma-separated values) is the most common data exchange format. Pandas handles it with \texttt{read\_csv}, which has over 50 parameters for dealing with real-world quirks.

\begin{minted}{python}
import pandas as pd

# Basic load
df = pd.read_csv("titanic.csv")
print(df.shape)
print(df.head())
\end{minted}

\subsection{Common CSV problems}

\begin{minted}{python}
# Problem 1: Wrong delimiter
# Some European CSVs use semicolons
df_euro = pd.read_csv("data_eu.csv", sep=";", decimal=",")

# Problem 2: Encoding issues
df_fr = pd.read_csv("french_data.csv", encoding="latin-1")

# Problem 3: Messy headers
df = pd.read_csv("messy.csv", header=0, skiprows=[1, 2])
df.columns = df.columns.str.strip().str.lower().str.replace(" ", "_")
print(df.columns.tolist())
\end{minted}

\begin{preprocesstip}
Always standardize column names immediately after loading: lowercase, underscores, no spaces. This prevents subtle bugs from inconsistent capitalization.
\end{preprocesstip}

\subsection{Optimizing CSV loads}

\begin{minted}{python}
# Load only needed columns (saves memory)
cols = ["PassengerId", "Survived", "Pclass", "Sex", "Age", "Fare"]
df = pd.read_csv("titanic.csv", usecols=cols)

# Specify dtypes to save memory
dtypes = {"Survived": "int8", "Pclass": "int8", "Sex": "category"}
df = pd.read_csv("titanic.csv", usecols=cols, dtype=dtypes)
print(df.memory_usage(deep=True))

# Load in chunks for large files
chunks = pd.read_csv("large_file.csv", chunksize=50_000)
result = pd.concat([chunk[chunk["Age"] > 30] for chunk in chunks])
\end{minted}

\begin{warningbox}
Never load a multi-gigabyte CSV into memory all at once. Use \texttt{chunksize}, or switch to Parquet/Feather format. A 2\,GB CSV typically becomes 400\,MB in Parquet with faster load times.
\end{warningbox}

% ============================================================
\section{Loading Excel files}

\begin{minted}{python}
# Basic Excel load (requires openpyxl)
df = pd.read_excel("survey_results.xlsx", sheet_name="Sheet1")

# Load all sheets into a dictionary
all_sheets = pd.read_excel("multi_sheet.xlsx", sheet_name=None)
for name, sheet_df in all_sheets.items():
    print(f"Sheet '{name}': {sheet_df.shape}")

# Skip header rows and footers
df = pd.read_excel("report.xlsx", skiprows=3, skipfooter=2,
                    engine="openpyxl")
\end{minted}

\begin{datatip}
Excel files preserve formatting but lose type precision. Dates may load as strings, percentages as floats between 0 and 1, and currency as plain numbers. Always verify dtypes after loading Excel.
\end{datatip}

% ============================================================
\section{Loading JSON data}

\begin{minted}{python}
import json

# Flat JSON
df = pd.read_json("flat_data.json")

# Nested JSON (common from APIs)
with open("nested_data.json") as f:
    raw = json.load(f)

# Use json_normalize to flatten
from pandas import json_normalize

df = json_normalize(raw["results"],
                    record_path="measurements",
                    meta=["station_id", "station_name"])
print(df.head())
\end{minted}

\begin{minted}{python}
# Real example: loading Gapminder-style JSON
import requests

url = "https://api.worldbank.org/v2/country/BEN/indicator/SP.POP.TOTL?format=json&per_page=60"
response = requests.get(url)
data = response.json()

# The actual data is in data[1]
df_wb = pd.DataFrame(data[1])
print(df_wb[["date", "value"]].head(10))
\end{minted}

% ============================================================
\section{Loading from SQL databases}

\begin{minted}{python}
from sqlalchemy import create_engine

# SQLite (file-based, no server needed)
engine = create_engine("sqlite:///local_data.db")

# Load entire table
df = pd.read_sql("SELECT * FROM patients", engine)

# Load with filtering (push computation to the database)
query = """
SELECT patient_id, age, diagnosis, lab_value
FROM patients
WHERE age > 40
  AND diagnosis IS NOT NULL
ORDER BY lab_value DESC
LIMIT 10000
"""
df = pd.read_sql(query, engine)
print(df.shape)
\end{minted}

\begin{preprocesstip}
When loading from SQL, always push as much filtering as possible into the query. Loading 10 million rows into pandas and then filtering is much slower than letting the database filter first.
\end{preprocesstip}

% ============================================================
\section{Loading from APIs}

\begin{minted}{python}
import requests
import pandas as pd

# REST API example: Open Meteo weather API
url = "https://archive-api.open-meteo.com/v1/archive"
params = {
    "latitude": 6.36,    # Cotonou
    "longitude": 2.43,
    "start_date": "2024-01-01",
    "end_date": "2024-12-31",
    "daily": "temperature_2m_mean,precipitation_sum",
}

response = requests.get(url, params=params)
data = response.json()

df_weather = pd.DataFrame({
    "date": data["daily"]["time"],
    "temp_mean": data["daily"]["temperature_2m_mean"],
    "precip_mm": data["daily"]["precipitation_sum"],
})
df_weather["date"] = pd.to_datetime(df_weather["date"])
print(df_weather.head())
\end{minted}

\begin{minted}{python}
# Pagination: many APIs return data in pages
all_records = []
page = 1

while True:
    resp = requests.get(f"{base_url}?page={page}&per_page=100")
    records = resp.json()["results"]
    if not records:
        break
    all_records.extend(records)
    page += 1

df = pd.DataFrame(all_records)
print(f"Total records: {len(df)}")
\end{minted}

% ============================================================
\section{Web scraping basics}

\begin{minted}{python}
import requests
from bs4 import BeautifulSoup
import pandas as pd

# Scrape an HTML table
url = "https://en.wikipedia.org/wiki/List_of_countries_by_GDP_(nominal)"
response = requests.get(url)
soup = BeautifulSoup(response.content, "html.parser")

# pandas can read HTML tables directly
tables = pd.read_html(url)
print(f"Found {len(tables)} tables on the page")
gdp_df = tables[2]  # pick the right table by index
print(gdp_df.head())
\end{minted}

\begin{warningbox}
Web scraping may violate a website's terms of service. Always check \texttt{robots.txt} and the site's ToS before scraping. For research, prefer official APIs and open data portals.
\end{warningbox}

% ============================================================
\section{Comparing formats: performance}

\begin{minted}{python}
import time

# Compare CSV vs Parquet load times
# First, create a Parquet version
df = pd.read_csv("large_dataset.csv")
df.to_parquet("large_dataset.parquet")

# Time CSV load
start = time.time()
df_csv = pd.read_csv("large_dataset.csv")
csv_time = time.time() - start

# Time Parquet load
start = time.time()
df_pq = pd.read_parquet("large_dataset.parquet")
pq_time = time.time() - start

print(f"CSV load:     {csv_time:.2f}s")
print(f"Parquet load: {pq_time:.2f}s")
print(f"Speedup:      {csv_time / pq_time:.1f}x")
\end{minted}

% ============================================================
\section{Exercises}

\begin{exercise}
\begin{enumerate}
  \item Load the Titanic CSV with optimized dtypes. Compare memory usage before and after dtype optimization.
  \item Load data from the World Bank API for three African countries (Benin, Nigeria, Ghana) for the indicator \texttt{SP.POP.TOTL} (total population). Combine into a single DataFrame.
  \item Load an Excel file with multiple sheets. Write a function that returns a dictionary mapping sheet names to their shapes.
  \item Scrape an HTML table from Wikipedia and clean the column names (lowercase, underscores, no special characters).
  \item Convert the Adult Census UCI dataset from CSV to Parquet. Compare file sizes and load times.
\end{enumerate}
\end{exercise}

% ============================================================
\section{Chapter summary}

\begin{itemize}
  \item Pandas provides readers for CSV, Excel, JSON, SQL, Parquet, HTML, and more.
  \item CSV loading requires attention to delimiters, encodings, headers, and dtypes.
  \item JSON from APIs often needs \texttt{json\_normalize} to flatten nested structures.
  \item SQL loads should push filtering into the query, not into pandas.
  \item APIs require handling pagination, rate limits, and authentication.
  \item Parquet is 3--5x faster than CSV for loading and uses 5x less disk space.
\end{itemize}
